\section{Evaluation}
\label{sec:eval}

We evaluate \Name through large-scale simulation and a physical RoCEv2 testbed, organized around six questions: \textbf{EQ1}~(\S\ref{sec:eq1}): under static workloads with scarce bandwidth, which schedulers actually keep premium-tier SLOs? \textbf{EQ2}~(\S\ref{sec:eq2}): when the workload is adversarially constructed to mislead static priority keys, how large is the penalty? \textbf{EQ3}~(\S\ref{sec:eq3}): when tenant demand changes mid-run, who tracks it? \textbf{EQ4}~(\S\ref{sec:eq4}): where does the gap come from mechanistically? \textbf{EQ5}~(\S\ref{sec:eq5}): do the simulation findings hold under controlled multi-job contention on physical hardware? \textbf{EQ6}~(\S\ref{sec:eq6}): does the enforcement path work with real NCCL training jobs?

\subsection{Experimental Setup}
\label{sec:setup}

\subsubsection{Simulation Methodology}
\label{sec:sim-method}

We use a custom simulator that explicitly models the per-iteration structure of distributed training: each job alternates between a computation phase with no network demand and a collective communication phase whose duration depends on its allocated bandwidth, so phase staggering (\S\ref{sec:periodic-origin}) emerges from the model rather than being assumed. The simulator implements the exact SLO semantics of \S\ref{sec:model}: a job's SLO is satisfied when its window-averaged iteration time $\bar{A}_i$ stays within a factor $c_i$ of its contention-free iteration time $T_i^{\mathrm{solo}}$, with anchor quantities ($T_i^{\mathrm{comp}}$, $T_i^{\mathrm{comm,solo}}$, and per-iteration byte volume) drawn from profiled model configurations. We validate simulator fidelity against the physical testbed by replicating static testbed scenarios bit-for-bit in configuration: the simulator reproduces the testbed's measured allocations and attainment within $0.2\%$, giving us confidence that simulation findings transfer to hardware (further cross-validation in \S\ref{sec:eq5}).

\subsubsection{Workloads and Scenarios}
\label{sec:workloads}

Table~\ref{tab:workloads} summarizes the five frozen scenario families. \textbf{E1} (feasibility ladder) sweeps the shared-spine capacity across the feasibility boundary of a fixed 14-job workload ($\Sigma_{\mathcal{P}}b^{\mathrm{att}}=605.7$\,Gbps, $\Sigma_{\mathcal{S}}b^{\mathrm{att}}=263.6$\,Gbps, frontier $C^{*}=737.5$\,Gbps), exposing scheduler behavior from deep infeasibility (400\,Gbps) to abundance (1200\,Gbps). \textbf{E2}$'$ (adversarial static) constructs a tenant mix in which the premium tier consists of large models whose static GPU-intensity ranking is the \emph{lowest} in the cluster---maximally misleading for intensity-keyed schedulers---while \textbf{E2-pro} is the orthogonal control in which intensity ranking and SLO tightness are positively correlated. \textbf{E3}/\textbf{E3}$'$ (dynamic) start from the E1 workload configuration and execute a mid-run regime swap at $t{=}300$\,s that reassigns slack coefficients, inverting which tenants hold tight SLOs: E3 ($800$\,Gbps) is the control arm where the swap is benign, and E3$'$ ($630$\,Gbps) is the kill arm in which the post-swap premium tier is exactly the set of jobs that static keys deprioritize (post-swap $\Sigma'_{\mathcal{P}}=569.6$\,Gbps, $\Sigma'_{\mathcal{S}}=230.0$\,Gbps, $C^{*\prime}=684.6$\,Gbps). All scenario definitions, including per-job model types, parallelism degrees, slack coefficients, and anchor constants, were frozen before any evaluation run and are versioned with every result artifact.

\begin{table}[t]
\centering
\caption{Frozen scenario families used throughout the evaluation.}
\label{tab:workloads}
\scriptsize
\setlength{\tabcolsep}{3pt}
\begin{tabular}{@{}l l c c p{2.2cm}@{}}
\toprule
\textbf{Scenario} & \textbf{Type} & \textbf{Jobs} & \textbf{Spine (Gbps)} & \textbf{Key Property} \\
\midrule
E1  & static ladder & 14 (8P/6S) & 400--1200 & feasibility ladder across $C^{*}{=}737.5$ \\
E2$'$ & adversarial static & 14 (10P/4S) & 400--800 & tier-model mapping inverted vs.\ static keys \\
E2-pro & favorable static & 14 (10P/4S) & 630, 800 & tier-model alignment matches static keys \\
E3  & dynamic control & 14 (8P/6S) & 800 & benign mid-run tier swap \\
E3$'$ & dynamic kill & 14 (8P/6S) & 630 & post-swap premium set deprioritized by static keys \\
\bottomrule
\end{tabular}
\end{table}


\subsubsection{Baselines}
\label{sec:baselines}

We compare against four baselines spanning the design space of \S\ref{sec:related-comparison}: \textbf{Fair}, work-conserving per-flow fair sharing; \textbf{SP}, static flow-size-based priority; \textbf{CRUX}~\cite{Crux}, static GPU-intensity-keyed priority; and \textbf{DF}, a windowed deficit-feedback controller in the family of \S\ref{sec:why-closed-form}---cross-class deficit round-robin with intra-class weights proportional to $\exp(K\pi_i)$ ($K{=}2$, clipping ratio $10{\times}$)---which we implement with its SLO targets derived from the same anchor semantics as \Name, so that no baseline loses on definitional grounds.

\subsubsection{Metrics and Statistical Methodology}
\label{sec:metrics}

Our primary metric is \emph{premium SLO attainment} (P-attn): the fraction of premium-tier jobs whose SLO attainment ratio---$\mathrm{SAS}_i = c_i T_i^{\mathrm{solo}}/\bar{A}_i$---meets or exceeds $0.98$ within the measurement window, a $2\%$ tolerance that absorbs window-quantization noise in window averaging.
Secondary metrics are \emph{P-cap} (the distribution of premium attainment ratios, capturing partial credit), \emph{S-cont} (standard-tier service level, detecting winners that sacrifice the floor), and \emph{starved jobs} (jobs receiving near-zero useful throughput). Static experiments run $600$\,s; dynamic experiments evaluate three pre-registered windows (W1 pre-swap, W2 immediately post-swap, W3 steady-state post-swap), and trajectory figures use sliding $100$\,s windows with regime-correct SLO targets. Every simulation number is the mean $\pm$ standard deviation over five seeds (in the E3 family, DF trajectories use seeds $\{0,1,2,3,4\}$ and all other policies $\{0,1,2,4,5\}$); the anchor experiment (Table~\ref{tab:anchor}) and the physical testbed campaign (\S\ref{sec:eq5}) follow separate protocols disclosed in their respective captions.

\subsubsection{Physical Testbed}
\label{sec:testbed-setup}

Our testbed consists of two GPU workers connected through a commodity switch at 100\,Gbps RoCEv2, running real NCCL training jobs instrumented with \Name's shim (\S\ref{sec:design}). We construct controlled contention scenarios in which a tight-SLO job competes with co-located traffic, and measure \emph{slowdown}: the ratio of the job's window time under contention to its solo window time. Each scenario repeats across independent rounds with alternating scheduler order to cancel drift; we report per-round values rather than pooled averages so that statistical ties are visible. The testbed also serves two forensic purposes: probing the full DSCP range to expose the traffic-class mapping behavior of commodity NICs, and quantifying NIC classification-granularity constraints (\S\ref{sec:impl}).

% ============================================================
% EQ1 — Bandwidth-Scarcity Ladder
% ============================================================
\subsection{Bandwidth-Scarcity Ladder}
\label{sec:eq1}

Does \Name differentiate from baselines under scarce bandwidth while remaining cost-free when bandwidth is abundant? To answer this, we sweep the spine bandwidth of a 16-node FatTree from 400 to 1200~Gbps using the frozen 14-job workload of \S\ref{sec:workloads} (8 premium, 6 standard) and measure P-attn at each capacity point ($n{=}5$ seeds, mean $\pm$ population std). Figure~\ref{fig:e1-ladder} reports per-policy P-attn across the ladder, with shaded regions marking scarce ($\leq500$~Gbps), transition (500--630~Gbps), and abundant ($\geq800$~Gbps) regimes. At 400~Gbps---the most severely constrained regime---\Name delivers $75.0\%$ P-attn ($\pm15.8\%$) against a best baseline of $62.5\%$ (CRUX), a gap of $12.5$ percentage points (pp); at 500~Gbps, \Name reaches $87.5\%$ while the best baseline (CRUX/SP) stalls at $70.0\%$ (a $17.5$~pp gap); at 630~Gbps, \Name attains $97.5\%$ against a best baseline of $87.5\%$ (Fair), a $10.0$~pp margin. Starting at 800~Gbps, all policies converge above $90.0\%$ (by 1000~Gbps \Name, DF, CRUX, and SP all reach $97.5\%$) and the differentiation disappears, consistent with the capacity-saturation expectation: the closed-form allocator does not over-engineer the abundant case. The baseline anchor at 400~Gbps (Table~\ref{tab:anchor}) confirms that this differentiation is not an artifact of workload imbalance: under a separate 24-job anchor workload, CRUX achieves the highest Mean SAS among baselines ($0.8422$, $29.2\%$ SLO attainment), yet all baselines leave substantial SLO slack that \Name is designed to eliminate. \textbf{Takeaway:} \Name delivers clear differentiation in the scarce and transition regimes---with P-attn margins of $10.0$--$17.5$~pp over the best baseline---and incurs no cost in the abundant regime, where all policies converge.

\begin{figure}[t]
\centering
\includegraphics[width=\columnwidth]{figure/fig2_e1_ladder.pdf}
\caption{E1 bandwidth ladder: P-attn vs.\ spine bandwidth for 14 jobs (8 premium, 6 standard), $n=5$ seeds, mean $\pm$ population std. Shaded regions denote scarce ($\leq500$~Gbps), transition (500--630~Gbps), and abundant ($\geq800$~Gbps) regimes.}
\label{fig:e1-ladder}
\end{figure}

\begin{table}[t]
\centering
\caption{Baseline anchor at 400 Gbps (24-job anchor workload, 12L/8M/4S, 3 seeds; anchor experiments use an independent 3-seed protocol distinct from the 5-seed standard of \S\ref{sec:metrics}).}
\label{tab:anchor}
\small\begin{tabular}{lcccc}\toprule
Policy & Mean SAS & SLO attainment & Collapse (\%) \\\midrule
LongLiu & 1.0000 & 100.0 & 0.0 \\
Fair & 0.7835 & 23.6 & 1.4 \\
CRUX & 0.8422 & 29.2 & 13.9 \\
SP & 0.7776 & 29.2 & 37.5 \\
DF & 0.8131 & 22.2 & 0.0 \\
\bottomrule\end{tabular}\end{table}

% ============================================================
% EQ2 — Orthogonal Workload Structure
% ============================================================
\subsection{Orthogonal Workload Structure}
\label{sec:eq2}

Is \Name's advantage specific to CRUX's favorable workload structure, or is it a structural property of its mechanism? To disentangle workload-driven effects from mechanism-driven effects, we construct two orthogonal workload variants: E2$'$, where the tier-to-model mapping is inverted to disadvantage CRUX's static key (the same 14-job structure with reversed premium/standard assignment), and E2-pro, a CRUX-favorable positive control that preserves the original tier alignment. Both use a spine-bandwidth ladder ($n{=}5$ seeds) and are reported in Figure~\ref{fig:e2-killfield} and Table~\ref{tab:e2pro}. In the disadvantaging E2$'$ scenario at 500~Gbps, \Name achieves $84.5\%$ P-attn ($\pm8.9\%$) against CRUX's $64.4\%$ ($\pm17.8\%$), a $20.1$~pp gap; at 630~Gbps, \Name reaches $91.1\%$ ($\pm13.0\%$) while \emph{all} baselines remain at or below $77.8\%$ (Fair, CRUX, SP at $77.8\%$; DF at $68.9\%$), a lead of at least $13.3$~pp. The gap persists because CRUX's static tier-to-key mapping---optimized for the original workload---cannot recover when the mapping is inverted, whereas \Name's allocation is recomputed from online anchors and is workload-agnostic. At 800~Gbps the gap closes as capacity saturates (\Name and DF both at $97.8\%$). In the E2-pro positive control, \Name, DF, and CRUX all achieve $100.0\%$ P-attn at both 630 and 800~Gbps, with SP and Fair close behind at $96.0\%\pm8.0\%$, confirming that when the workload structure is favorable, all policies perform adequately. \textbf{Takeaway:} \Name's advantage is structural, not an artifact of workload alignment with CRUX's static key; the $20.1$~pp gap under the inverted mapping shows that \Name is robust to workload-structure mismatch, while CRUX collapses when its favorable alignment is removed.

\begin{figure}[t]
\centering
\includegraphics[width=\columnwidth]{figure/fig3_e2_orthogonal.pdf}
\caption{E2 orthogonal workload: (a)~E2$'$ disadvantaging CRUX with inverted tier-to-model mapping, and (b)~E2-pro favorable CRUX positive control. Both use $n=5$ seeds, mean $\pm$ population std. A single shared legend applies to both panels.}
\label{fig:e2-killfield}
\end{figure}

\begin{table}[t]
\centering
\caption{E2-pro positive control (CRUX-favorable, 5 seeds).}
\label{tab:e2pro}
\small\begin{tabular}{lcc}\toprule
Policy & 630 Gbps & 800 Gbps \\\midrule
LongLiu & 100.0$\pm$0.0\% & 100.0$\pm$0.0\% \\
DF & 100.0$\pm$0.0\% & 100.0$\pm$0.0\% \\
CRUX & 100.0$\pm$0.0\% & 100.0$\pm$0.0\% \\
SP & 96.0$\pm$8.0\% & 96.0$\pm$8.0\% \\
Fair & 96.0$\pm$8.0\% & 96.0$\pm$8.0\% \\
\bottomrule\end{tabular}\end{table}

% ============================================================
% EQ3 — Dynamic Tier Reconfiguration
% ============================================================
\subsection{Dynamic Tier Reconfiguration}
\label{sec:eq3}

Can \Name maintain SLO guarantees when the tier assignment of jobs changes mid-run---a scenario that stresses any priority scheduler's ability to reconfigure without transient degradation? We test this with a two-arm swap experiment at $t{=}300$~s ($n{=}5$ seeds): the control arm (E3, 800~Gbps) swaps tier assignments in a direction that advantages CRUX's static mapping, while the kill arm (E3$'$, 630~Gbps) swaps in the opposite direction, disadvantaging CRUX (Figure~\ref{fig:eval-e3}). In both arms, \Name maintains $100.0\pm0.0\%$ P-attn across all three windows (W1 pre-swap, W2 transient, W3 post-swap), with the sliding windows truncated at the regime boundary ($t{=}300$~s) so that no cross-regime mixing occurs. CRUX reaches $63.3\pm16.3\%$ in the control arm and collapses to $10.0\pm12.2\%$ in the kill arm, while DF achieves $96.7\pm6.7\%$ (control) and $25.0\pm22.4\%$ (kill). The key mechanism is \Name's closed-form allocation: per-job bandwidth targets are recomputed instantly from the new tier assignment with no iterative convergence delay, so the system transits from one feasible allocation to another without transient SLO violations. \textbf{Takeaway:} \Name provides instantaneous, transient-free adaptation to tier reconfiguration, maintaining $100.0\pm0.0\%$ P-attn throughout the swap in both arms, while CRUX and DF suffer severe post-swap degradation ($10.0\%$ and $25.0\%$ respectively) in the kill arm.

\begin{figure}[t]
\centering
\includegraphics[width=\columnwidth]{figure/fig1_hero.pdf}
\caption{Dynamic tier swap: sliding-window P-attn trajectories for \Name, DF, and CRUX in (a)~E3 control arm at 800~Gbps and (b)~E3$'$ kill arm at 630~Gbps, $n=5$ seeds, mean $\pm$ population std. Windows are truncated at the regime boundary ($t=300$~s), so no cross-regime mixing occurs. Gray shaded regions denote W1 (pre-swap), W2 (transient), and W3 (post-swap) windows; the dashed vertical line marks the tier-swap event. A single shared figure legend applies to both panels.}
\label{fig:eval-e3}
\end{figure}

% ============================================================
% EQ4 — Why the Deficit-Feedback Controller Fails
% ============================================================
\subsection{Why the Deficit-Feedback Controller Fails}
\label{sec:eq4}

Why does DF, despite achieving high pre-swap P-attn in the E3 control arm, fail to converge correctly after the tier swap in the kill arm? We examine DF's sliding-window P-attn trajectory in both arms (Figure~\ref{fig:d1-transient}) and the per-job priority ($\pi$) time series in E3$'$ (Figure~\ref{fig:d1-mechanism}). In the E3 control arm, DF transitions from $70.0\%$ (W1) through a transient dip at the swap boundary to $96.7\%$ (W3), converging toward the correct post-swap premium set. In the E3$'$ kill arm, DF starts at $48.0\%$ (W1), dips to $36.0\%$ at the swap, briefly overshoots to a transient peak of $62.5\%$ (t$\,\approx\,$399~s), and then \emph{declines} to $30.0\%$ at W3---converging in the wrong direction, toward starvation of the post-swap premium jobs it should be protecting. (The fixed-window metric of \S\ref{sec:eq3} reads $25.0\%$ at W3, differing from the sliding-window endpoint by half a window width.) The trajectory exposes both a cold-start transient (W1 ramp-up) and a swap-induced transient (W2 overshoot followed by W3 degradation).

The $\pi$ time series (Figure~\ref{fig:d1-mechanism}, E3$'$ DF, seed~0) reveals the mechanism. Post-swap, the eight new-premium jobs (J5--J12) settle into a heterogeneous $\pi$ stall range of $[+0.056, +0.399]$: LLM jobs cluster at $\pi \approx +0.35$--$+0.40$, while J6 (LLaMA-2-13B) stalls at $\pi \approx +0.06$ and J11 (T5-11B) at $+0.26$. This $7\times$ spread in $\pi$ within the premium tier drives a $\sim 1.9\times$ spread in $\exp(\pi \cdot K)$ scheduling weights, preventing uniform premium protection: premium jobs with lower $\pi$ are under-fed while well-fed jobs over-consume. The same evidence rules out the three alternative explanations. It is not capacity exhaustion: the weight-clipping mechanism (\texttt{clip\_ratio}=10) never triggers, the maximum within-tier weight ratio ($\sim 2\times$) lying far below the threshold. It is not competition loss: DF still allocates $\sim 78$--$80\%$ of aggregate scheduling weight to the premium pool, comparable to \Name's dedicated premium provisioning, yet P-attn falls to $30.0\%$. And the standard tier is neither starved nor overfed in reverse: the old-premium-now-standard jobs are compressed to $\sim 20$--$22\%$ of the total pool.
The root cause is an \emph{expressiveness wall}: DF's scalar $\pi$ cannot encode tier isolation and per-job capping simultaneously, the two properties \Name's two-level closed form expresses directly. In contrast, \Name exhibits zero variance across all seeds and windows ($\text{std} = 0.0$ vs.\ DF's $\pm 22.4$~pp in the kill arm W3). \textbf{Takeaway:} DF's feedback law is executed with machine precision yet fails structurally, because a single-layer $\pi$ space lacks the representational capacity of a per-job target vector---an expressiveness wall that \Name's closed-form design eliminates, together with the transient dynamics that accompany any feedback law.

\begin{figure}[t]
\centering
\includegraphics[width=\columnwidth]{figure/fig4_d1_trajectory.pdf}
\caption{DF sliding-window P-attn trajectory in (a)~E3 control arm at 800~Gbps and (b)~E3$'$ kill arm at 630~Gbps, $n=5$ seeds, mean $\pm$ population std. Annotated values (70.0\%, 96.7\%, 48.0\%, 30.0\%) are fixed-window means at W1 ($t\approx250$~s) and W3 ($t\approx550$~s) midpoints. Gray shaded regions denote W1--W3 windows; the dashed line marks the tier swap. The \Name reference line (100\%) is shown for comparison. A single shared legend applies to both panels.}
\label{fig:d1-transient}
\end{figure}

\begin{figure}[t]
\centering
\includegraphics[width=\columnwidth]{figure/fig5_pi_timeseries.pdf}
\caption{Per-job $\pi$ time series for DF in the E3$'$ kill arm (seed~0): (a)~jobs that are premium before the swap, and (b)~jobs that are standard before the swap. Each job is identified by its ID and model abbreviation (e.g., ``J0 (13B)''), with unique marker--linestyle pairs for redundancy coding. The vertical dashed line marks the tier swap; the horizontal dashed line at $\pi=0$ separates over-fed ($\pi<0$) from under-fed ($\pi>0$) regimes. A single external shared legend covers all 14 jobs.}
\label{fig:d1-mechanism}
\end{figure}

% ============================================================
% EQ5 — CPU Emulator Ladder (Multi-Job Validation)
% ============================================================
\subsection{CPU Emulator Ladder}
\label{sec:eq5}

Do the simulation findings hold under controlled multi-job contention on physical hardware? We validate using a CPU-based RDMA emulator that isolates the scheduler logic from GPU-specific dynamics. The emulator runs on a two-node RoCEv2 testbed (same hardware as \S\ref{sec:testbed-setup}) and implements the exact SLO semantics of the simulator: each job alternates between computation (no network demand) and communication phases, with per-job DSCP marking under scheduler control. We compare \Name against two baselines---Static (fixed priority: premium at P6, standard at P2) and Fair (no priority differentiation)---across two scenarios: \textbf{S1} (5 jobs, 2 premium + 3 standard) with four regime depths (ample, moderate, deep, very deep), and \textbf{S2} (6 jobs, 3 premium + 3 standard) representing starvation conditions. Each regime-arm combination repeats for 5 independent rounds with alternating scheduler order to cancel drift.

Figure~\ref{fig:expC-s1-pattn} reports the P-attn metric (sum of premium slowdown excess, lower is better) across S1 regimes. In the moderate regime, \Name achieves $0.28$ P-attn against Static's $0.44$, a $36\%$ improvement; in the deep regime, \Name reaches $0.57$ vs.\ Static's $0.62$ (an $8\%$ improvement). The ample regime shows convergence (all policies near $0.3$--$0.4$), confirming that \Name incurs no cost when bandwidth is abundant. The very deep regime exhibits Fair outperforming both \Name and Static ($0.22$ vs.\ $0.27$/$0.30$), attributed to phase mutual exclusion effects under extreme contention where fair-sharing becomes optimal. This pattern---\Name outperforming Static in scarcity, matching in abundance---mirrors the simulation prediction of \S\ref{sec:eq1}.

Figure~\ref{fig:expC-s2-perjob} presents per-job slowdown in the S2 starvation scenario. For premium jobs J0 and J1 (high communication volume), all three arms fail to meet the SLO target ($c_{eval}=1.5$), with slowdown ranging from $2.27$ (Fair) to $3.00$ (Static/LongLiu)---a structural impossibility when three 8\,MB jobs contend for a 50\,Gbps link. The key differentiator is J2 (lower communication volume): \Name reduces J2's slowdown to $1.43$ against Static's $1.97$ (a $28\%$ improvement), demonstrating dynamic priority protection for the marginally-solvable premium job. For standard jobs (J3--J5), Static achieves near-solo performance ($\approx 0.99\times$) due to phase staggering, while \Name's standard jobs show bounded degradation ($\le 1.16\times$), consistent with Theorem~1's worst-case guarantee. Fair exhibits uniform premium-standard starvation ($\approx 1.3\times$ across all jobs), validating the ``Rule 2'' design constraint that fair-sharing violates SLOs under structural overload.

\textbf{Takeaway:} On controlled multi-job hardware, \Name delivers $8$--$36\%$ P-attn improvement over Static in scarcity regimes (moderate/deep), converges in abundance, and exhibits bounded standard degradation in starvation conditions---consistent with simulation predictions and theory bounds. The key limitation is the mlx5 NIC's lack of per-priority QoS, which attenuates DSCP differentiation effects; this is a hardware constraint, not a scheduler design flaw.

\begin{figure}[t]
\centering
\includegraphics[width=\columnwidth]{figure/fig_expC_s1_pattn.pdf}
\caption{S1 P-attn across regimes for 5-job workload (2 premium + 3 standard). Lower is better. LongLiu outperforms Static in moderate and deep regimes by 36\% and 8\%, respectively, and converges in the ample regime. Fair outperforms in very deep regime due to phase mutual exclusion. Error bars: $\pm$1 std across 5 rounds.}
\label{fig:expC-s1-pattn}
\end{figure}

\begin{figure}[t]
\centering
\includegraphics[width=\columnwidth]{figure/fig_expC_s2_perjob.pdf}
\caption{S2 per-job slowdown for 6-job starvation scenario (3 premium + 3 standard). Horizontal dashed line: SLO target $c_{eval}=1.5$. For structurally-infeasible premium jobs (J0/J1), all arms fail. For marginally-solvable J2, LongLiu achieves $1.43\times$ vs.\ Static's $1.97\times$ (28\% improvement). Static standard jobs show near-solo performance ($\approx 0.99\times$); LongLiu standard degradation bounded at $\le 1.16\times$. Error bars: $\pm$1 std across 5 rounds.}
\label{fig:expC-s2-perjob}
\end{figure}

% ============================================================
% EQ6 — Physical Testbed Validation (NCCL)
% ============================================================
\subsection{Physical Testbed Validation (NCCL)}
\label{sec:eq6}

Does \Name's dynamic priority adjustment translate from simulation to a real RDMA fabric under contention? We answer this on a two-node single-bottleneck dumbbell testbed (Fig.~\ref{fig:testbed-v6}, Table~\ref{tab:testbed-hw}) with a 6~Gbps iperf3 UDP background flow at DSCP~P3 under strict-priority queueing, $c_i=1.2$ for tight jobs, and 5-minute warmup before each arm.

\begin{table}[t]
\centering
\caption{Physical testbed hardware.}
\label{tab:testbed-hw}
\scriptsize
\setlength{\tabcolsep}{3pt}
\begin{tabular}{@{}l l p{4cm}@{}}
\toprule
\textbf{Component} & \textbf{Model} & \textbf{Role} \\
\midrule
GPU & Quadro RTX 4000 (8\,GB) & classifying NIC (fw 22.41.1000) \\
GPU& Quadro RTX 5000 (16\,GB) & non-classifying NIC (fw 32.41.1000) \\
NIC & 2$\times$BlueField-3 B3220 (100\,Gbps) & classifying vs.\ single queue \\
Switch & commodity 100\,GbE & DSCP-to-CoS, strict-priority \\
Link   & RoCEv2 (50\,Gbps) & single bottleneck dumbbell \\
\bottomrule
\end{tabular}
\end{table}

On this single bottleneck, the closed-form allocation of \S\ref{sec:design} degenerates to a demand check: unmet premium jobs receive strict priority, satisfied jobs maintain best-effort handling.
We compare \Name (dynamic $P6{\rightarrow}P2$) against CRUX (static P3) across four independent rounds (two orderings LL$\rightarrow$CX and CX$\rightarrow$LL $\times$ two replications; a third replication, rep1, was verified via md5sum to be a byte-identical copy and excluded), each comprising two staggered jobs of 15 windows---phase~1 (solo and ramp-up, windows~0--6) followed by phase~2 (dual-job contention, windows~7--14).

\paragraph{Stable contention window (windows 7--11)}
Over the primary stable-contention window, \Name achieves a mean slowdown of $1.11\times$ against CRUX at $1.26$--$1.36\times$ (Table~\ref{tab:eq5-summary}), a relative advantage of $13.7$--$16.8\%$. Three of the four independent rounds (orig\_r1, orig\_r2, rep2\_r2) exhibit non-overlapping slowdown intervals---the tightest case showing \Name $[1.08, 1.14]$ vs.\ CRUX $[1.26, 1.30]$ at orig\_r2. The remaining round (rep2\_r1, LL$\rightarrow$CX ordering) produces overlapping intervals (\Name $1.25\times \pm 0.08$ vs.\ CRUX $1.25\times \pm 0.01$), attributed to an anomalous \Name elevation that preserves the qualitative trend. Considering all four rounds gives a replication rate of 3/4 decisive non-overlap with 1/4 statistically indistinguishable.

\paragraph{Full window (windows 7--14, sensitivity)}
The full window, which includes end-of-run contention decay as the competing job completes, serves as a sensitivity check: \Name remains at $1.09$--$1.21\times$ while CRUX spans $1.10$--$1.34\times$, with the CRUX low end ($1.10\times$ at orig\_r1) an artifact of the competing job's completion rather than steady-state contention. The stable-window metric remains the primary comparison point.

\paragraph{Deployment forensics}
Two findings from this testbed generalize beyond \Name and are detailed in \S\ref{sec:impl}: a silent DSCP-to-traffic-class inversion, in which the default $\texttt{trafficClass} = \mathtt{priority} \times 8$ rule maps the highest software label to the lowest hardware class---invisible to connectivity checks and surfaced only by a full-class probe verified against NIC \texttt{tx\_prio} counters---and a NIC-generation classification asymmetry, in which one testbed NIC funnels all DSCP values into a single hardware priority queue, sinking the enforcement surface to the first switch hop. Both are properties of specific hardware generations, not of the scheduler design, but any marking-based deployment must probe for them rather than assume correct mapping.

\textbf{\Name reduces tight-SLO job slowdown by $13.7$--$16.8\%$ relative to the strongest baseline on commodity hardware---three rounds decisive with non-overlapping intervals, one statistically tied---consistent with the simulation prediction of \S\ref{sec:eq2}.}

\begin{figure*}[t]
\centering
\includegraphics[width=\linewidth]{figure/fig6_testbed.pdf}
\caption{Physical testbed comparison of \Name{} (dynamic, blue solid line) and CRUX (static~P3, red dashed line) across four independent rounds. Subfigure titles indicate run ordering (\Name$\rightarrow$CRUX or CRUX$\rightarrow$\Name). The gray band marks the stable-contention window (windows~7--11), the primary metric; windows~7--14 are given as a sensitivity reference. rep2\_r1 (bottom-left) shows an anomalous \Name{} elevation producing overlapping intervals (1/4 overlap). Round~rep1 excluded via md5 byte-identity verification. $n{=}1$ (single run per round; replications are independent experiment repetitions, not seeds).}
\label{fig:testbed-v6}
\end{figure*}

\begin{table}[t]
\centering
\caption{Stable contention window (windows~7--11): per-round mean and [minimum, maximum] slowdown for \Name{} and CRUX. Rows marked with $\dagger$ show non-overlapping intervals. Rep1 excluded via md5 byte-identity verification.}
\label{tab:eq5-summary}
\footnotesize
\begin{tabular}{lcccc}
\toprule
\textbf{Round} & \textbf{Order} & \textbf{Scheduler} & \textbf{Mean} & \textbf{[Min, Max]}\\
\midrule
\multirow{2}{*}{orig\_r1${}^\dagger$}
   & \multirow{2}{*}{\Name~$\rightarrow$~CRUX}
   & \Name  & 1.1075 & {[}1.0780, 1.1441{]} \\
   &                & CRUX  & 1.3022 & {[}1.2942, 1.3101{]} \\
\multirow{2}{*}{orig\_r2${}^\dagger$}
   & \multirow{2}{*}{CRUX~$\rightarrow$~\Name}
   & \Name  & 1.1069 & {[}1.0840, 1.1346{]} \\
   &                & CRUX  & 1.2824 & {[}1.2586, 1.3034{]} \\
\multirow{2}{*}{rep2\_r1}
   & \multirow{2}{*}{\Name~$\rightarrow$~CRUX}
   & \Name  & 1.2496 & {[}1.1378, 1.2909{]} \\
   &                & CRUX  & 1.2541 & {[}1.2478, 1.2618{]} \\
\multirow{2}{*}{rep2\_r2${}^\dagger$}
   & \multirow{2}{*}{CRUX~$\rightarrow$~\Name}
   & \Name  & 1.1123 & {[}1.0978, 1.1482{]} \\
   &                & CRUX  & 1.3365 & {[}1.2785, 1.3825{]} \\
\bottomrule
\end{tabular}
\end{table}

% ============================================================
% Summary
% ============================================================
\subsection{Summary}
\label{sec:eval-summary}

Across five questions the evidence converges on a single pattern. In static regimes, \Name stratifies the field exactly where bandwidth is scarce and costs nothing where it is abundant (EQ1); when the workload is built to mislead static keys, \Name widens its lead while keyed schedulers collapse, and it matches the best baseline when keys happen to be favorable (EQ2). In dynamic regimes---the operating condition of real clusters---\Name is the only scheduler that sustains every premium SLO through an instantaneous tier reversal, while both the static-key and the deficit-feedback designs fail, each for a structural reason that no parameter repair addresses (EQ3, EQ4). The physical testbed confirms that these findings are not simulation artifacts: the enforcement path works on unmodified commodity hardware, and the two pitfalls it surfaces are documented with their detection method and fix (EQ5). Together, the five answers instantiate the three contributions of this paper: the feasibility frontier predicts where schedulers must diverge, the closed-form allocator is what survives those regimes, and the hardware enforcement layer carries the allocation onto real wire.